\documentclass[journal]{vgtc}                     % final (journal style)
%\documentclass[journal,hideappendix]{vgtc}        % final (journal style) without appendices
%\documentclass[review,journal]{vgtc}              % review (journal style)
%\documentclass[review,journal,hideappendix]{vgtc} % review (journal style)
%\documentclass[widereview]{vgtc}                  % wide-spaced review
%\documentclass[preprint,journal]{vgtc}            % preprint (journal style)


%% Uncomment one of the lines above depending on where your paper is
%% in the conference process. ``review'' and ``widereview'' are for review
%% submission, ``preprint'' is for pre-publication in an open access repository,
%% and the final version doesn't use a specific qualifier.

%% If you are submitting a paper to a conference for review with a double
%% blind reviewing process, please use one of the ``review'' options and replace the value ``0'' below with your
%% OnlineID. Otherwise, you may safely leave it at ``0''.
\onlineid{0}

%% In preprint mode you may define your own headline. If not, the default IEEE copyright message will appear in preprint mode.
%\preprinttext{To appear in IEEE Transactions on Visualization and Computer Graphics.}

%% In preprint mode, this adds a link to the version of the paper on IEEEXplore
%% Uncomment this line when you produce a preprint version of the article
%% after the article receives a DOI for the paper from IEEE
%\ieeedoi{xx.xxxx/TVCG.201x.xxxxxxx}

%% declare the category of your paper, only shown in review mode
\vgtccategory{Research}

%% Paper title.
\title{Uso de modelo de Machine Learning para detección de genes de resistencia a antibióticos en microbiota intestinal porcina}

%% Author ORCID IDs should be specified using \authororcid like below inside
%% of the \author command. ORCID IDs can be registered at https://orcid.org/.
%% Include only the 16-digit dashed ID.
\author{ Juan Jerónimo Gómez Rubiano }

\authorfooter{
  %% insert punctuation at end of each item
  \item
  	Universidad Nacional de Colombia, Facultad de Ingeniería
}

%% Abstract section.
\abstract{%
  \lipsum[1] % filler text. Replace with your abstract.
  %
  %% We recommend that you link to your supplemental material here in the abstract, as well
  %% as in the Supplemental Materials section at the end.
  A free copy of this paper and all supplemental materials are available at \url{https://OSF.IO/2NBSG}.
}

%% Keywords that describe your work. Will show as 'Index Terms' in journal
%% please capitalize first letter and insert punctuation after last keyword
\keywords{Radiosity, global illumination, constant time}

%% A teaser figure can be included as follows
\teaser{
  \centering
  \includegraphics[width=\linewidth, alt={A view of clouds with orange sunrays shining through from behind.}]{clouds}
  \caption{%
  	Dramatic evening clouds.
  	Note that the teaser may not be wider than the abstract block.%
  }
  \label{fig:teaser}
}

%% Uncomment below to disable the manuscript note
%\renewcommand{\manuscriptnotetxt}{}

%% Copyright space is enabled by default as required by guidelines.
%% It is disabled by the 'review' option or via the following command:
%\nocopyrightspace


%%%%%%%%%%%%%%%%%%%%%%%%%%%%%%%%%%%%%%%%%%%%%%%%%%%%%%%%%%%%%%%%
%%%%%%%%%%%%%%%%%%%%%% LOAD PACKAGES %%%%%%%%%%%%%%%%%%%%%%%%%%%
%%%%%%%%%%%%%%%%%%%%%%%%%%%%%%%%%%%%%%%%%%%%%%%%%%%%%%%%%%%%%%%%

%% Tell graphicx where to find files for figures when calling \includegraphics.
%% Note that due to the \DeclareGraphicsExtensions{} call it is no longer necessary
%% to provide the the path and extension of a graphics file:
%% \includegraphics{diamondrule} is completely sufficient.
\graphicspath{{figs/}{figures/}{pictures/}{images/}{./}} % where to search for the images

%% Only used in the template examples. You can remove these lines.
\usepackage{booktabs}                  % only used for the table example
\usepackage{lipsum}                    % used to generate placeholder text
\usepackage{mwe}                       % used to generate placeholder figures
\usepackage{ccicons}                   % package to be able to use icons from creative commons

%% We encourage the use of mathptmx for consistent usage of times font
%% throughout the proceedings. However, if you encounter conflicts
%% with other math-related packages, you may want to disable it.
\usepackage{mathptmx}                  % use matching math font

\begin{document}

%%%%%%%%%%%%%%%%%%%%%%%%%%%%%%%%%%%%%%%%%%%%%%%%%%%%%%%%%%%%%%%%
%%%%%%%%%%%%%%%%%%%%%% START OF THE PAPER %%%%%%%%%%%%%%%%%%%%%%
%%%%%%%%%%%%%%%%%%%%%%%%%%%%%%%%%%%%%%%%%%%%%%%%%%%%%%%%%%%%%%%%

%% The ``\maketitle'' command must be the first command after the
%% ``\begin{document}'' command. It prepares and prints the title block.
%% the only exception to this rule is the \firstsection command
\firstsection{Introducción}

\maketitle

La resistencia por parte de los distintos patógenos a los antibióticos, es uno de los mayores problemas de salud pública en la actualidad. Es tanto un problema de escala global en el presente, como una amenaza emergente para el futuro. En los últimos años (2019-2025) ha aumentado el número de muertes asociadas a microorganismos resistentes de 1.9 millones a 4 millones anualmente, con los últimos pronósticos, advirtiendo que esta cifra va a aumentar a los 10 millones de muertes anuales si no se logran medidas de gran impacto para evitarlo. 

Existen varios factores biológicos y sociales que conllevan a que distintas especies y cepas de microorganismos desarrollen resistencia. De manera natural, todas las especies de microorganismos desarrollan mecanismos bioquímicos para resistir a antibióticos, y las cepas que desarrollan estos mecanismos se vuelven cada vez más prevalentes por selección natural. Es decir, realmente no hay una manera de detener este fenómeno de manera permanente. 

No obstante, si es posible que se acelere el desarrollo de la resistencia, y la forma más frecuente de hacerlo, es por medio del suministro de antibióticos de manera preventiva o terapéutica. Evidentemente, sí es necesario utilizar antibióticos para curar enfermedades, porque muchas enfermedades infecciosas son letales sin tratamiento. Sin embargo, el uso inadecuado de antibióticos, particularmente, el uso excesivo de estos, naturalmente acelera el desarrollo de mecanismos de resistencia, al punto que los tratamientos que antes contenían una determinada cepa, ya no son eficaces. 

El principal lugar en donde esto ha sido reportado como un problema son las infecciones nosocomiales (estadía hospitalaria), ya que, irónicamente, es donde más se utiizan antibióticos de alto espectro. Las bacterias multirresistentes tradicionalmente son menos virulentes, como un tipo de "trade-off" evolucionario que se ha documentado por un buen tiempo (entre más resistente, menos patogénica), pero ciertamente son capaces de causar infecciones, especialmente en pacientes inmunosuprimidos, o vulnerables por otros factores (edad, desnutrición, etc.). El reporte global de Naghavi et al., muestra que en los últimos años la mortalidad asociada a resistencia antibiótica ha disminuido significativamente en la población pediátrica, pero ha aumentado en los adultos mayores. Un análisis más profundo corrobora las estancias hospitalarias como factor de riesgo de morbimortalidad por microorganismos multirresistentes.

Si bien los infectólogos suelen tener una "escalera" de antibióticos, y definen una serie de antibióticos específicos para cepas más resistentes, varias cepas eventualmente evolucionan al punto de que son resistentes a todos los antibióticos conocidos en el planeta. O al menos, son suficientemente resistentes para causar grandes sobrecostos en el uso de medicamentos, o causar la muerte en el paciente antes de que se pueda aplicar una terapia eficaz.

Entre los factores asociados con el desarrollo de resistencia a los antibióticos, uno de los factores de mayor impacto, pero también más subestimados, ha sido el uso de antibióticos en los sectores de la ganadería y cría de animales para consumo humano. Los antibióticos son usados de manera industrial y sistemática en las granjas por dos grandes motivos: Primero, de manera preventiva, ya que, al prevenir el desarrollo de infecciones en los animales, se evita la mortalidad prematura, y la infestación de especies patógenas en la carne. Por otro lado, los antibióticos actúan como factores de crecimiento, lo que evidentemente, hace que los animales alcancen más rápido la madurez, aumentando la productividad y ganancias respectivas de la granja.

Se estima que aproximadamente 93.000 toneladas de antibióticos son usadas cada año con fines preventivos en granjas y criaderos a nivel mundial. Un estudio de Li et al. (2025), advierte que el resistoma de las granjas es más diverso que el resistoma humano, o el de suelo. Además, identifican los resistomas como particularmente diversos entre especies, notando que las medidas y análisis en cerdos debe ser especializado del de pollos o el de vacas por ejemplo. Esto es particularmente notorio porque el uso indiscriminado de anitbióticos en ganado se está convirtiendo en un nuevo problema emergente de salud pública, adicional al uso incorrecto de antibióticos en el sector médico. En particular, varias cepas que ya son patógenas en humanos, como la CC398 de \textit{Staphylococcus aureus} Metilicino-Resistente, tienen plásmidos de resistencia que se originaron en granjas, y pudieron transmitirse a los humanos. Es un peligro en salud pública, tanto a nivel alienticio, como para los trabajadores del sector ganadero.

La Inteligencia Artificial se ha presentado como una de las herramientas novedosas más prometedoras para hacer frente a este problema. El mecanismo más claro en el que puede ayudar es en el desarrollo de nuevos antibióticos. La halicina fue el primer fármaco descubierto por Inteligencia Artificial, y sirve para tratar cepas de bacterias como SARM o \textit{Acinetobacter baumanii} que antes de su descubrimiento eran resistentes a todos los antibióticos conocidos. Si bien su farmacocinética y toxicidad todavía no están lo suficientemente investigadas para lanzarlo al mercado, este descubrimiento abrió la puerta a la posibilidad de usar modelos de IA para descubrir nuevos antibióticos. Desde entonces, ... entre otros han surgido como nuevos fármacos en la lucha contra la resistencia antimicrobiana, particularmente de última línea.

No obstante, no se puede dejar todo el trabajo al descubrimiento optimizado de fármacos. Un aspecto crucial del uso de IA en este contexto, es la detección temprana de cepas multirresistentes. Ya se utiliza IA para predecir entornos en los cuales pueden surgir microorganismos resistentes antes de que tengan impacto directo en la salud pública, lo que es crucial para mitigar sus impactos, ya sea tanto 

\section{Metodología}

\subsection{Fuente de datos y selección de muestras}

Se utilizaron datos metagenómicos crudos del proyecto PRJEB11755 del European Nucleotide Archive (ENA), un estudio de la microbiota intestinal porcina organizado en siete grupos de granja/cohorte (identificados en los metadatos como PIG, EYZ, SYZ, BMZ, BHZ, ZXZ y DB), con 200, 52, 38, 34, 24, 22 y 20 corridas de secuenciación profundas disponibles por grupo, respectivamente. Cada muestra biológica cuenta en ENA con dos corridas asociadas: una profunda (del orden de 2--5\,GB comprimidos) y una acompañante de tamaño muy inferior ($<$100\,MB), insuficiente para un ensamblado de novo confiable; se descartaron sistemáticamente estas últimas filtrando por tamaño.

Dado que este trabajo está diseñado para ejecutarse en un computador portátil comercial, se seleccionó un subconjunto de 50 corridas, mantenido proporcional al tamaño relativo de cada grupo (26/7/5/4/3/3/2 corridas para PIG/EYZ/SYZ/BMZ/BHZ/ZXZ/DB, respectivamente, garantizando al menos una corrida por grupo), en vez de usar el total de más de 390 corridas disponibles. Adicionalmente, ya que procesar una corrida completa (hasta 47 millones de pares de lecturas) es computacionalmente costoso para un portátil, se aplicó un submuestreo aleatorio reproducible de 1\,000\,000 de pares de lecturas por muestra (\texttt{seqtk}, semilla fija) previo al ensamblado. Debido a interrupciones recurrentes durante la descarga remota desde los servidores de ENA --- caídas de conexión y bloqueos temporales tipo HTTP 403, atribuibles a limitación de tasa del servidor --- al momento de este análisis se contaba con 32 de las 50 corridas seleccionadas completamente procesadas; los resultados de este artículo corresponden a ese subconjunto.

\subsection{Preprocesamiento y detección de ARG por comparación con catálogo}

Para cada corrida se aplicó el siguiente flujo, ejecutado de forma independiente e idempotente por muestra: control de calidad de las lecturas (\texttt{fastp}), submuestreo reproducible (\texttt{seqtk}), ensamblado de novo (\texttt{MEGAHIT}), predicción de genes/ORFs sobre los contigs ensamblados (\texttt{Prodigal}), y detección de genes de resistencia antimicrobiana (ARG) por homología de secuencia contra la base de datos curada CARD, usando la herramienta RGI (BLASTP/DIAMOND) en modo local. Este último paso constituye el método de referencia tipo "diccionario": únicamente detecta secuencias suficientemente similares a un modelo ya catalogado en CARD, y es el resultado que se busca contrastar con el modelo de aprendizaje automático descrito a continuación.

\subsection{Conjunto de datos para el modelo de aprendizaje automático}
\label{sec:dataset_ml}

Se plantea la pregunta de si un modelo puede aprender un patrón composicional asociado a "gen de resistencia" a partir únicamente de las secuencias de referencia de CARD, y si dicho patrón generaliza a secuencias reales de las muestras, incluyendo ORFs que RGI nunca marcó. Para ello se construyeron dos fuentes de datos deliberadamente separadas:

\begin{itemize}
  \item \textbf{Positivos de entrenamiento.} Se extrajeron las secuencias de ADN de referencia de las 6.404 entradas de CARD, agrupadas en 470 familias de genes según su categoría ARO. Como las secuencias de CARD corresponden a genes completos (longitud mediana $\approx$876\,pb), mientras que los ORFs reales son fragmentos de un ensamblaje metagenómico submuestreado (longitud mediana $\approx$360\,pb en nuestras muestras), cada secuencia de CARD se complementó con tres recortes aleatorios de longitud muestreada de la distribución real observada. Esto evita que el modelo aprenda a distinguir "secuencia completa" en vez de un patrón de resistencia genuino, y produjo 25.616 secuencias positivas de entrenamiento.
  \item \textbf{Negativos de entrenamiento y evaluación del mundo real.} De los 358.543 ORFs reales predichos por Prodigal en las 32 muestras, los 453 que RGI ya identificó como ARG se apartaron como conjunto de evaluación del mundo real y nunca se usaron para entrenar; los 358.090 ORFs restantes constituyen el \emph{pool} negativo de entrenamiento (muestreado 3:1 respecto a los positivos).
\end{itemize}

\subsection{Extracción de características y modelos}

Se usó como representación la frecuencia relativa de tetranucleótidos (k-mers de ADN, $k=4$, 256 dimensiones), una firma composicional libre de alineamiento, comúnmente usada en tareas de \emph{binning} taxonómico. Sobre esta representación se entrenaron y compararon dos modelos de clasificación binaria (ARG vs.\ no-ARG): regresión logística con balanceo de clases, como línea base, y \texttt{XGBoost} (300 árboles, profundidad máxima 5), como modelo principal.

\subsection{Esquemas de evaluación}

Se evaluaron dos particiones sobre los datos derivados de CARD (\cref{sec:dataset_ml}):

\begin{itemize}
  \item \textbf{Partición aleatoria}: reparto 80/20 estratificado por clase; sirve como verificación de cordura, pero sobreestima la generalización real, ya que fragmentos de una misma familia de gen pueden quedar repartidos entre entrenamiento y prueba.
  \item \textbf{Partición por familia de gen (\emph{holdout})}: se reservaron 94 de las 470 familias completas de CARD exclusivamente para prueba, de forma que ninguna variante de esas familias es vista en entrenamiento. Esta es la prueba de generalización genuina: mide si el modelo reconoce algo composicionalmente nuevo, en vez de memorizar secuencias parecidas a las de entrenamiento.
\end{itemize}

Finalmente, se entrenó un modelo con la totalidad de los datos derivados de CARD y se aplicó sobre los 358.543 ORFs reales de las 32 muestras, comparando sus predicciones contra las detecciones de RGI mediante dos indicadores: (1) el \emph{recall} sobre los 453 ARG ya conocidos, y (2) el número de ORFs no catalogados por RGI que el modelo marca como ARG probable ("candidatos nuevos"). Un subconjunto de los candidatos de mayor probabilidad se verificó puntualmente mediante BLAST remoto (\emph{megablast}) contra la base \texttt{nt} de NCBI.

\section{Resultados}

\subsection{Muestras procesadas y datos disponibles}

Al momento de este análisis, 32 de las 50 corridas seleccionadas de PRJEB11755 se procesaron exitosamente de principio a fin (descarga, ensamblado, predicción de genes y detección de ARG); las 18 restantes se vieron afectadas por interrupciones de red o bloqueos temporales del servidor de ENA durante la descarga, y quedan pendientes para una corrida posterior del pipeline. Sobre estas 32 muestras, Prodigal predijo 358.543 ORFs, de los cuales RGI confirmó 453 como genes de resistencia por homología contra CARD (0,13\,\% del total), proporción consistente con lo esperado en una muestra metagenómica no enriquecida para ARG.

\subsection{Desempeño del modelo sobre los datos de CARD}

La \cref{tab:model_results} resume el desempeño de ambos modelos en las dos particiones de evaluación.

\begin{table}[tb]
  \caption{Desempeño de los modelos de clasificación (positivo = ARG) en las dos particiones de evaluación sobre los datos derivados de CARD.}
  \label{tab:model_results}
  \small
  \centering
  \begin{tabular}{llccc}
    \toprule
    Partición & Modelo & ROC-AUC & AUPRC & Recall \\
    \midrule
    Aleatoria & Reg.\ logística & 0.922 & 0.844 & 0.812 \\
    Aleatoria & XGBoost         & 0.982 & 0.967 & 0.902 \\
    Por familia (\emph{holdout}) & Reg.\ logística & 0.830 & 0.623 & 0.576 \\
    Por familia (\emph{holdout}) & XGBoost         & 0.930 & 0.840 & 0.636 \\
    \bottomrule
  \end{tabular}
\end{table}

En la partición aleatoria, ambos modelos alcanzan un desempeño alto (XGBoost: ROC-AUC=0,982, AUPRC=0,967), aunque, como se anticipó, esta partición sobreestima la capacidad de generalización real. La partición por familia --- la prueba más exigente, al retener 94 de las 470 familias completas fuera del entrenamiento --- muestra que XGBoost mantiene un desempeño considerablemente por encima del azar (ROC-AUC=0,930, AUPRC=0,840, \emph{recall}=0,636 sobre la clase positiva), mientras que la regresión logística cae de forma más marcada (ROC-AUC=0,830, AUPRC=0,623, \emph{recall}=0,576). Esto sugiere que las interacciones no lineales entre frecuencias de k-mers que captura XGBoost son relevantes para generalizar a familias de genes nunca vistas, y que el patrón composicional de "gen de resistencia" tiene, en efecto, una señal real más allá de la memorización de secuencias específicas.

\subsection{Desempeño sobre datos reales y candidatos nuevos}

Al aplicar el modelo final (entrenado con la totalidad del conjunto derivado de CARD) sobre los 358.543 ORFs reales, el \emph{recall} sobre los 453 ARG ya confirmados por RGI fue de apenas 37,5\,\%, sustancialmente menor que el \emph{recall} observado en la partición por familia sobre datos de CARD (63,6\,\%). Esta caída sugiere una brecha de dominio entre las secuencias de referencia de CARD (genes completos, sin errores de secuenciación ni de ensamblado) y los ORFs reales derivados de un ensamblaje metagenómico submuestreado (fragmentos posiblemente truncados en el borde de un contig, con errores propios del ensamblado y la secuenciación).

El modelo marcó adicionalmente 9.828 ORFs (2,7\,\% del \emph{pool} negativo) como ARG probable sin que RGI los hubiera catalogado. Dado el \emph{recall} limitado sobre los positivos ya confirmados, es esperable que la precisión real sobre este conjunto de "candidatos nuevos" sea baja. Como verificación puntual, se sometieron a BLAST remoto (\emph{megablast} contra \texttt{nt}) los cinco candidatos de mayor probabilidad ($\geq$0,98); ninguno correspondió a un gen de resistencia conocido: uno tuvo como mejor coincidencia un genoma ensamblado de un bacteriófago (\textit{Caudoviricetes} sp.), y los cuatro restantes coincidieron con genomas bacterianos completos (\textit{Alistipes} sp., \textit{Salmonella enterica}, y \textit{Streptococcus alactolyticus} en dos muestras distintas) sin anotación funcional de resistencia asociada. Esto es consistente con falsos positivos originados en similitud composicional genérica (p.\,ej., contenido GC o sesgo de codones típico de ciertos genomas bacterianos), más que en homología funcional real con un gen de resistencia.

\section{Conclusiones}

\subsection{Loading figures}

The style automatically looks for image files with the correct extension (eps for regular \LaTeX; pdf, png, and jpg for pdf\LaTeX), in a set of given subfolders defined above using \verb|\graphicspath|: figures/, pictures/, images/.
It is thus sufficient to use \verb|\includegraphics{clouds}| (instead of \verb|\includegraphics{figs/clouds.jpg}|).
Figures should be in CMYK or Grey scale format, otherwise, colour shifting may occur during the printing process.

\subsection{Vector figures}

Vector graphics like \texttt{pdf} and \texttt{eps} (can be generated from a \texttt{svg}) are best for charts and other figures with text or lines.
They will look much nicer and crisper and any text in them will be more selectable, searchable, and accessible.

\subsection{Raster figures}

Of the raster graphics formats, screenshots of user interfaces and text, as well as line art, are better shown with \texttt{png}.
\texttt{jpg} is better for photographs.
Make sure all raster graphics are captured in high enough resolution so they look crisp and scale well.

\subsection{Alt texts}

Add alternative texts that describe the content of the image to all figures.

\subsection{Figures on the first page}

The teaser figure should only have the width of the abstract as the template enforces it.
The use of figures other than the optional teaser is not permitted on the first page.
Other figures should begin on the second page.
Papers submitted with figures other than the optional teaser on the first page will be refused.

\subsection{Subfigures}

You can add subfigures using the \texttt{subcaption} package that is automatically loaded.
Inside a \verb|figure| environment, create a \verb|subfigure| environment.
See \cref{fig:ex_subfigs} for an example.
You can reference individual figures, either fully using \verb|\cref| (\cref{fig:ex_subfigs_a,fig:ex_subfigs_b}) or by letter using \verb|\subref|.
E.g., \subref{fig:ex_subfigs_b}, \subref{fig:ex_subfigs_c}.
Note that \verb|\subref| only works for one label at a time.

\begin{figure}[tbp]
  \centering
  \begin{subfigure}[b]{0.45\columnwidth}
  	\centering
  	\includegraphics[width=\textwidth, alt={Big letter A on a gray background.}]{example-image-a}
  	\caption{The letter A.}
  	\label{fig:ex_subfigs_a}
  \end{subfigure}%
  \hfill%
  \begin{subfigure}[b]{0.45\columnwidth}
  	\centering
  	\includegraphics[width=\textwidth, alt={Big letter B on a gray background.}]{example-image-b}
  	\caption{The letter B.}
  	\label{fig:ex_subfigs_b}
  \end{subfigure}%
  \\%
  \begin{subfigure}[b]{0.45\columnwidth}
  	\centering
  	\includegraphics[width=\textwidth, alt={Big letter C on a gray background.}]{example-image-c}
  	\caption{The letter C.}
  	\label{fig:ex_subfigs_c}
  \end{subfigure}%
  \subfigsCaption{Example of adding subfigures with the \texttt{subcaption} package.}
  \label{fig:ex_subfigs}
\end{figure}

\subsection{Figure Credits}
\label{sec:figure_credits_inst}

In the \hyperref[sec:figure_credits]{Figure Credits} section at the end of the paper, you should credit the original sources of any figures that were reproduced or modified.
Include any license details necessary, as well as links to the original materials whenever possible.
For credits to figures from academic papers, include a citation that is listed in the \textbf{References} section.
An example is provided \hyperref[sec:figure_credits]{below}.


\section{Equations and Tables}

Equations can be added like so:

\begin{equation}
  \label{eq:sum}
  \sum_{j=1}^{z} j = \frac{z(z+1)}{2}
\end{equation}

Tables, such as \cref{tab:vis_papers} can also be included.

\begin{table}[tb]
  \caption{%
  	VIS/VisWeek accepted/presented papers: 1990--2025, data from \href{https://www.vispubdata.org/}{vispubdata} \cite{Isenberg2017Vispubdata.orgMetadataCollection}.%
  }
  \label{tab:vis_papers}
  \scriptsize%
	\newlength{\digitwidth}%
  \settowidth{\digitwidth}{0}%
  \centering%
  \begin{tabular}{%
  	  r%
  	  	*{8}{@{\hspace{5.5pt}}c@{\hspace{5.5pt}}}%
  	  	*{2}{@{\hspace{5.5pt}}r}%
  	}
  	\toprule
  	year & \rotatebox{90}{VIS (from 2021)} & \rotatebox{90}{Vis/SciVis} &   \rotatebox{90}{SciVis conf} &   \rotatebox{90}{InfoVis} &   \rotatebox{90}{VAST} &   \rotatebox{90}{VAST conf} &   \rotatebox{90}{TVCG @ VIS} &   \rotatebox{90}{CG\&A @ VIS} &   \rotatebox{90}{VIS/VisWeek} \rotatebox{90}{incl.\ TVCG/CG\&A}   &   \rotatebox{90}{VIS/VisWeek} \rotatebox{90}{w/o TVCG/CG\&A}   \\
  	\midrule
  	2025 & 131 &    &   &    &    &    & 75 & 12 & 218 & 131 \\
  	2024 & 124 &    &   &    &    &    & 65 & 12 & 201 & 124 \\
  	2023 & 133 &    &   &    &    &    & 60 & \hspace{\digitwidth}9 & 202 & 133 \\
  	2022 & 119 &    &   &    &    &    & 67 & 18 & 204 & 119 \\
  	2021 & 109 &    &   &    &    &    & 50 & 11 & 170 & 109 \\
  	2020 &     & 32 &   & 64 & 51 & 10 & 41 & 12 & 210 & 157 \\
  	2019 &     & 25 &   & 53 & 42 & \hspace{\digitwidth}9 & 47 & 12 & 188 & 129 \\
  	2018 &     & 32 &   & 47 & 41 & \hspace{\digitwidth}7 & 42 & 12 & 181 & 127 \\
  	2017 &     & 23 &   & 39 & 37 & 15 & 21 & \hspace{\digitwidth}8 & 143 & 114 \\
  	2016 &     & 30 &   & 37 & 33 & 15 & 23 & 10 & 148 & 115 \\
  	2015 &     & 33 & 9 & 38 & 33 & 14 & 17 & 15 & 159 & 127 \\
  	2014 &     & 34 &   & 45 & 33 & 21 & 20 &    & 153 & 133 \\
  	2013 &     & 31 &   & 38 & 32 &    & 20 &    & 121 & 101 \\
  	2012 &     & 42 &   & 44 & 30 &    & 23 &    & 139 & 116 \\
  	2011 &     & 49 &   & 44 & 26 &    & 20 &    & 139 & 119 \\
  	2010 &     & 48 &   & 35 & 26 &    &    &    & 109 & 109 \\
  	2009 &     & 54 &   & 37 & 26 &    &    &    & 117 & 117 \\
  	2008 &     & 50 &   & 28 & 21 &    &    &    &  99 &  99 \\
  	2007 &     & 56 &   & 27 & 24 &    &    &    & 107 & 107 \\
  	2006 &     & 63 &   & 24 & 26 &    &    &    & 113 & 113 \\
  	2005 &     & 88 &   & 31 &    &    &    &    & 119 & 119 \\
  	2004 &     & 70 &   & 27 &    &    &    &    &  97 &  97 \\
  	2003 &     & 74 &   & 29 &    &    &    &    & 103 & 103 \\
  	2002 &     & 78 &   & 23 &    &    &    &    & 101 & 101 \\
  	2001 &     & 74 &   & 22 &    &    &    &    &  96 &  96 \\
  	2000 &     & 73 &   & 20 &    &    &    &    &  93 &  93 \\
  	1999 &     & 69 &   & 19 &    &    &    &    &  88 &  88 \\
  	1998 &     & 72 &   & 18 &    &    &    &    &  90 &  90 \\
  	1997 &     & 72 &   & 16 &    &    &    &    &  88 &  88 \\
  	1996 &     & 65 &   & 12 &    &    &    &    &  77 &  77 \\
  	1995 &     & 56 &   & 18 &    &    &    &    &  74 &  74 \\
  	1994 &     & 53 &   &    &    &    &    &    &  53 &  53 \\
  	1993 &     & 55 &   &    &    &    &    &    &  55 &  55 \\
  	1992 &     & 53 &   &    &    &    &    &    &  53 &  53 \\
  	1991 &     & 50 &   &    &    &    &    &    &  50 &  50 \\
  	1990 &     & 53 &   &    &    &    &    &    &  53 &  53 \\
  	\midrule
  	\textbf{sum} & \textbf{616} & \textbf{1657} & \textbf{9} & \textbf{835} & \textbf{481} & \textbf{91} & \textbf{591} & \textbf{131} & \textbf{4411} & \textbf{3689} \\
  	\bottomrule
  \end{tabular}%
\end{table}

\begin{figure}[tb]% specify a combination of t, b, p, or h for top, bottom, on its own page, or here
  \centering % avoid the use of \begin{center}...\end{center} and use \centering instead (more compact)
  \includegraphics[width=\columnwidth, alt={A line graph showing paper counts between 0 and 160 from 1990 to 2016 for 9 venues.}]{paper-count-2024}
  \caption{%
  	A visualization of the 1990--2024 data from \cref{tab:vis_papers}, recreated based on Fig.\ 1 from \cite{Isenberg2017Vispubdata.orgMetadataCollection}.%
  }
  \label{fig:vis_papers}
\end{figure}


\section{Supplemental Material Instructions}
\label{sec:supplement_inst}

In support of transparent research practices and long-term open science goals, you are encouraged to make your supplemental materials available on a publicly-accessible repository.
Please describe the available supplemental materials in the \hyperref[sec:supplemental_materials]{Supplemental Materials} section.
These details could include (1) what materials are available, (2) where they are hosted, and (3) any necessary omissions.

\section{Reporting of User Studies}

Please note that for the reporting of any experimental results that involve human participants you are \textbf{required} to ``include a statement in the article that the research was performed under the oversight of an institutional review board or equivalent local/regional body, including the official name of the IRB/ethics committee, or include an explanation as to why such a review was not conducted. For research involving human subjects, authors shall also report that consent from the human subjects in the research was obtained or explain why consent was not obtained'' \cite[Section 8.1.1.E]{IEEEPublications2025}. Ideally, for an IRB approval or similar, you include the case number under which the permission was granted. For instance:
``Our experiment was approved by our university's IRB (No. 12345678). [...] At the start of the experiment, we obtained informed consent from all participants, who filled in and signed a consent form (which we share in our additional materials).''

\section{References}
\label{sec:references_inst}

An example of the reference formatting is provided in the \textbf{References} section at the end.


\subsection{Include DOIs}
\label{sec:references_inst:doi}

All references which have a DOI (which are virtually all entries of a list of references today), by default, should have it included in the bib\TeX\ file such that they are displayed in the list of references (as a courtesy for your reviewers and your readers).
The DOI can be entered with or without the \url{https://doi.org/} prefix. Note that you can also use short DOIs; for these use the interface from \href{https://shortdoi.org/}{\texttt{shortdoi\discretionary{}{.}{.}org}} to obtain a short version of any valid DOI, as showcased in the example \textbf{References} section at the end of the template.

\subsection{Narrow DOI option}

The \verb|-narrow| versions of the bibliography style use the font \verb|PTSansNarrow-TLF| for typesetting the DOIs in a compact way.
This font needs to be available on your \LaTeX\ system.
It is part of the \href{https://www.ctan.org/pkg/paratype}{\texttt{paratype} package}, and many distributions (such as MikTeX) have it automatically installed.
If you do not have this package yet and want to use a \verb|-narrow| bibliography style then use your \LaTeX\ system's package installer to add it.
If this is not possible you can also revert to the respective bibliography styles without the \verb|-narrow| in the file name.
DVI-based processes to compile the template apparently cannot handle the different font so, by default, the template file uses the \texttt{abbrv-doi} bibliography style.

\subsection{Disabling hyperlinks}

To avoid adding hyperlinks to the references (the default) you can use \verb|\bibliographystyle{abbrv-doi}| instead of \verb|\bibliographystyle{abbrv-doi-hyperref}|.
By default, the DOI field in a bib\TeX\ entry is turned into a hyperlink.

See the examples in the bib\TeX\ file and the bibliography at the end of this template.

\subsection{Guidelines for bibTeX}

\begin{itemize}
  \item All bibliographic entries should be sorted alphabetically by the last name of the first author.
        This \LaTeX/bib\TeX\ template takes care of this sorting automatically.
  \item Merge multiple references into one; e.\,g., use \cite{Max1995OpticalModelsDirect,Kitware2003VisualizationToolkitUsers} (not \cite{Kitware2003VisualizationToolkitUsers}\cite{Max1995OpticalModelsDirect}).
        Within each set of multiple references, the references should be sorted in ascending order.
        This \LaTeX/bib\TeX\ template takes care of both the merging and the sorting automatically.
  \item Verify all data obtained from digital libraries, even ACM's DL and IEEE Xplore  etc.\ are sometimes wrong or incomplete.
  \item Do not trust bibliographic data from other services such as Mendeley.com, Google Scholar, or similar; these are even more likely to be incorrect or incomplete.
  \item Articles in journal---items to include:
        \begin{itemize}
  	      \item author names
  	      \item title
  	      \item journal name
  	      \item year
  	      \item volume
  	      \item number
  	      \item (optionally) month of publication as variable name (i.e., \texttt{\{jan\}} for January, etc.; month ranges using \texttt{\{jan \#\{/\}\# feb\}} or \texttt{\{jan \#\{-{}-\}\# feb\}})
        \end{itemize}
  \item Use journal names in proper style: correct: ``\texttt{IEEE Transactions on Visualization and Computer Graphics}'', incorrect: ``\texttt{Visualization and Computer Graphics, IEEE Transactions on}'', also incorrect: ``\texttt{IEEE transactions on visualization and computer graphics}''.  You can also (consistently) use the ISO4 abbreviation standard for journal names. In this case TVCG would be ``\texttt{IEEE Trans Vis Comput Graph}''. You can search for ISO4 abbreviations at \href{https://journal-abbreviations.library.ubc.ca/}{\texttt{journal\discretionary{}{-}{-}abbreviations\discretionary{}{.}{.}library\discretionary{}{.}{.}ubc\discretionary{}{.}{.}ca}} or at \href{https://woodward.library.ubc.ca/woodward/research-help/journal-abbreviations/}{\texttt{woodward\discretionary{}{.}{.}library\discretionary{}{.}{.}ubc\discretionary{}{.}{.}ca\discretionary{/}{}{/}woodward\discretionary{/}{}{/}research\discretionary{}{-}{-}help\discretionary{/}{}{/}journal\discretionary{}{-}{-}abbreviations}} (if the full journal name is not available, search for its individual words).
  \item Papers in proceedings---items to include:
        \begin{itemize}
  	      \item author names
  	      \item title
  	      \item abbreviated proceedings name: e.g., ``\texttt{Proc.\textbackslash{} CONF\_ACRONYNM}'' \textbf{without the year}; example: ``\texttt{Proc.\textbackslash{} CHI}'', ``\texttt{Proc.\textbackslash{} 3DUI}'', ``\texttt{Proc.\textbackslash{} Eurographics}'', ``\texttt{Proc.\textbackslash{} EuroVis}''
  	      \item year
        \end{itemize}

  \item Article/paper title convention: refrain from using curly brackets, except for acronyms/proper names/words following dashes/question marks etc.; example:\\\\
        %
        The paper ``Marching Cubes: A High Resolution 3D Surface Construction Algorithm'' should be entered as ``\texttt{\{M\}arching \{C\}ubes: A High Resolution \{3D\} Surface Construction Algorithm}'' or  ``\texttt{\{M\}arching \{C\}ubes: A high resolution \{3D\} surface construction algorithm}''.
        It will then be typeset as ``Marching Cubes: A high resolution 3D surface construction algorithm''.
  \item For all entries:
        \begin{itemize}
  	      \item DOI (or short DOI, see \cref{sec:references_inst:doi}) can be entered in the DOI field as plain DOI number or as DOI url.
  	      \item ``\texttt{articleno}'' plus ``\texttt{numpages}'' or ``\texttt{pages}'': Provide full page ranges \texttt{AA-{}-BB}. Alternatively, if an article number is available like for recent ACM conference publications, use that plus the corresponding number of pages instead (e.g., see the entry for Panavas et al.\ \cite{Panavas2022JuvenileGraphicalPerception}).
        \end{itemize}
  \item When citing references, do not use the reference as a sentence object; e.g., wrong: ``In \cite{Lorensen1987MarchingCubesHigh} the authors describe \dots'', correct: ``Lorensen and Cline \cite{Lorensen1987MarchingCubesHigh} describe \dots''
\end{itemize}

\section{Appendices}
\label{sec:appendices_inst}

Appendices can be specified using \verb|\appendix|.
For example, our Troubleshooting instructions in
\iflabelexists{appendix:troubleshooting}
  {\cref{appendix:troubleshooting}}
  {the appendix of the full paper at \url{https://osf.io/nrmyc}}.

Note that the paper submission has to end after the \textbf{References} section and within the page limit of the conference you are submitting to.
Any version of Appendices or the paper with Appendices included has to be submitted separately as supplementary material (not that, for IEEE VIS starting in 2026, you can optionally also submit an additional full version including the appendices).
You can use the \verb|hideappendix| class option to remove everything after \verb|\appendix|.
We encourage you to submit a full version of your paper to a preprint server with any appendices included.

You can use the \verb|\iflabelexists| macro to cross reference an appendix from the main text, but only if that label (i.e., the appendix) actually exists.
For example, above we use

\begin{verbatim}
\iflabelexists{appendix:troubleshooting}
  {\cref{appendix:troubleshooting}}
  {the appendix of the full paper at
   \url{https://osf.io/XXXXX}}.
\end{verbatim}

In order to cross-reference to the appendix with \verb|\cref| if it exists, but if the appendix is commented out then we will simply create a hyperlinked URL to it.


\section{Filler Text to Flush Out the Paper}

\lipsum[1-2]% Just add some more arbitrary text so we see a fuller paper example


\section*{Supplemental Materials}
\label{sec:supplemental_materials}

Refer to the instructions for this section (\cref{sec:supplement_inst}).
Below is an example you can follow that includes the actual supplemental material for this template:

All supplemental materials are available on OSF at \href{https://doi.org/10.17605/OSF.IO/2NBSG}{\texttt{osf\discretionary{}{.}{.}io\discretionary{/}{}{/}2nbsg}}, released under a \href{https://creativecommons.org/licenses/by/4.0/}{\ccby{} CC BY 4.0 license}.
In particular, they include (1) Excel files containing the data for and analyses for creating \cref{tab:vis_papers} and \cref{fig:vis_papers}, (2) figure images in multiple formats, and (3) a full version of this paper with all appendices.
Our other code is intellectual property of a corporation---Starbucks Research---and there is no feasible way to share it publicly.


\section*{Figure Credits and Copyrights}
\label{sec:figure_credits}

Refer to the instructions for this section (\cref{sec:figure_credits_inst}).
Here are the actual figure credits for this template:

\Cref{fig:teaser} image credit: Dominik Moritz, licensed under \href{https://creativecommons.org/public-domain/cc0/}{CC0}.
\Cref{fig:vis_papers} is a partial recreation of Fig.\ 1 from \cite{Isenberg2017Vispubdata.orgMetadataCollection}. The new image is from \href{https://github.com/pisenberg/vispubdata/}{\texttt{github\discretionary{}{.}{.}com\discretionary{/}{}{/}pisenberg\discretionary{/}{}{/}vispubdata}} and we use it under a \href{https://creativecommons.org/licenses/by/4.0/}{\ccby{} CC BY 4.0 license}. All other figures (in particular, \cref{fig:ex_subfigs}) remain under the authors' own copyright, with the permission to be used here. We also share them under a \href{https://creativecommons.org/licenses/by/4.0/}{\ccby{} CC BY 4.0 license} at \href{https://doi.org/10.17605/OSF.IO/2NBSG}{\texttt{osf\discretionary{}{.}{.}io\discretionary{/}{}{/}2nbsg}}.


%% if specified like this the section will be omitted in review mode
\acknowledgments{%
	The authors wish to thank A, B, and C.
  This work was supported in part by a grant from XYZ (\# 12345-67890).%
}


\bibliographystyle{abbrv-doi-hyperref}
%\bibliographystyle{abbrv-doi-hyperref-narrow}
%\bibliographystyle{abbrv-doi}
%\bibliographystyle{abbrv-doi-narrow}

\bibliography{template}


\appendix % You can use the `hideappendix` class option to skip everything after \appendix
\crefalias{section}{appendix} % this is to make sure that cleverref switches to referring to Appx. X from here on

\section{About Appendices}
Refer to \cref{sec:appendices_inst} for instructions regarding appendices.

\section{Troubleshooting}
\label{appendix:troubleshooting}

\subsection{ifpdf error}

If you receive compilation errors along the lines of \texttt{Package ifpdf Error: Name clash, \textbackslash ifpdf is already defined} then please add a new line \verb|\let\ifpdf\relax| right after the \verb|\documentclass[journal]{vgtc}| call.
Note that your error is due to packages you use that define \verb|\ifpdf| which is obsolete (the result is that \verb|\ifpdf| is defined twice); these packages should be changed to use \verb|ifpdf| package instead.


\subsection{\texttt{pdfendlink} error}

Occasionally (for some \LaTeX\ distributions) this hyper-linked bib\TeX\ style may lead to \textbf{compilation errors} (\texttt{pdfendlink ended up in different nesting level ...}) if a reference entry is broken across two pages (due to a bug in \verb|hyperref|).
In this case, make sure you have the latest version of the \verb|hyperref| package (i.e.\ update your \LaTeX\ installation/packages).
If this package update does not fix the error for you, you can use the \verb|documentclass| option \verb|protectagainstpdfendlink| such as as follows:\\[1ex]
{\footnotesize\verb|\documentclass[journal]{vgtc,protectagainstpdfendlink}|}\\[1ex]
to set the \LaTeX/bib\TeX{} processing to avoid breaking reference entries in the list of references. Alternatively, you can revert back to \verb|\bibliographystyle{abbrv-doi}| (at the expense of removing hyperlinks from the bibliography) and then to try \verb|\bibliographystyle{abbrv-doi-hyperref}| again after some more editing.

\end{document}
